\section{Related Work}
\label{sec:related_work}

\subsection{Cellular KPI Forecasting and Time-Series Prediction}

Mobile and wireless networking research has studied prediction as a foundation for proactive management. Recent sixth-generation (6G) surveys emphasize programmable, data-driven radio access networks for telemetry-driven control and assurance \cite{wang2023road6g}. O-RAN studies connect radio access programmability with operational automation \cite{polese2023oran}. Recent RAN alarm prediction work reflects the value of modeling sparse alarms with network context \cite{li2025ranalarm}. Universal mobile-traffic forecasting similarly motivates learning cellular traffic dynamics \cite{chai2025uomo}.

General time-series forecasting now spans recurrent, convolutional, transformer, decomposition and language-model-inspired designs. TimesNet represents a recent temporal-variation modeling direction \cite{wu2023timesnet}, while strong linear baselines question whether transformer complexity is always necessary \cite{zeng2023transformers}. Follow-up models explore inverted transformer structures, multiscale mixing and probabilistic sequence pretraining \cite{liu2024itransformer,wang2024timemixer,ansari2024chronos}. Time-LLM and GPT4MTS adapt language-model ideas to time-series prediction \cite{jin2024timellm,jia2024gpt4mts}. Donut and OmniAnomaly further motivate learning from seasonal or multivariate KPI time series in production systems \cite{xu2018donut,su2019omni}. Cellular KPI forecasting studies mainly focus on traffic prediction or individual KPI modeling, whereas operational degradation prediction requires continuous KPI evolution, sparse fault events and cross-cell context. Existing approaches rarely address alarm-KPI coupling under enterprise monitoring conditions with incomplete topology information and expert RCA annotations. CET-STGL addresses this setting by integrating KPI forecasting, weak degradation estimation and alarm-candidate prioritization within an event-token representation framework.

\subsection{Alarm Correlation, RCA and Event Influence Modeling}

Alarm correlation is a classic fault-management problem in communication networks. Recent work has explored alarm-correlation rule mining for network fault management \cite{fournierViger2021alarmrules}, whereas log-based operational diagnosis methods such as DeepLog illustrate the broader role of event sequences in detecting and diagnosing operational anomalies \cite{du2017deeplog}.

Recent artificial intelligence for information technology operations (AIOps) and microservice RCA studies provide methodological background for multimodal observability and graph-based diagnosis. Nezha models multi-modal observability for interpretable root-cause analysis \cite{yu2023nezha}. CausalRCA and neural Granger causal discovery explore causal-inference-inspired localization \cite{xin2023causalrca,lin2024run}, and Chain-of-Event learns weighted event causal graphs for microservice incidents \cite{yao2024chainofevent}. RCAEval and AIOpsLab highlight the need for reproducible RCA evaluation infrastructure \cite{pham2025rcaeval,ma2025aiopslab}. Point-process analysis provides a foundation for self-exciting event sequences \cite{ogata1988pointprocess}, and Hawkes-based Granger modeling connects event intensity with temporal influence estimation \cite{xu2016grangerhawkes}. Existing alarm correlation and RCA methods often focus on post-incident localization. Proactive degradation prediction instead asks whether sparse alarms add useful context before KPI deterioration. In cellular operations, alarms are often noisy, delayed or incomplete, making direct event-based prediction challenging. CET-STGL treats alarms as contextual event signals and studies their contribution through event tokens, temporal intensity and influence modeling. Ranking results are interpreted as weak-label consistency rather than validated RCA, and event influence features are causal-inspired association priors rather than proof of causality.

\subsection{Spatio-Temporal Modeling and Telecom LLM Context}

Spatio-temporal graph neural networks model temporal dynamics and entity dependencies. Recent traffic-forecasting work has developed delay-aware transformers for long-range propagation effects \cite{jiang2023pdformer}. Adaptive spatio-temporal embeddings and large-scale traffic datasets support systematic evaluation of interacting entities \cite{liu2023staeformer,liu2023largest}. Recent KDD work on asynchronous spatio-temporal graph convolution also highlights irregular temporal structure in graph-based forecasting \cite{zhang2024aseer}. Existing spatio-temporal graph learning methods show the importance of interactions among network entities, but cellular degradation prediction remains limited by incomplete topology information and operational labels. CET-STGL adapts this idea to cellular monitoring by using inferred cross-cell context and train-only event influence priors. We evaluate a lightweight sklearn-based instantiation; end-to-end neural graph encoders are left for future evaluation under the same protocol.

Recent telecom and networking LLM work examines language models for standards knowledge, networking tasks and domain adaptation. TeleQnA evaluates whether LLMs capture domain-specific standards and networking concepts \cite{maatouk2026teleqna}. LLM-oriented telecom research helps position telecom-domain AI, but CET-STGL addresses a different data modality: structured operational telemetry for KPI forecasting, weak degradation classification and weak alarm-candidate prioritization on a private operations slice.
